 \documentclass[12pt]{article}[a4]
%\usepackage[utf8]{inputenc}
\usepackage[spanish,es-tabla]{babel}
\usepackage{amsmath,multicol}


\usepackage[backend=bibtex,sorting=none]{biblatex}
%usepackage[resetlabels,labeled]{multibib}
%\usepackage[sorting=none]{biblatex}

\addbibresource{referencias.bib}


\usepackage{fancyhdr}
\usepackage{import}
\usepackage{tikz}
\usepackage{booktabs}
\usepackage{graphicx,float,subcaption}
\newcommand{\gr}{^{\circ}}
\newcommand{\ohm}{\Omega}
\usepackage{import}
\usepackage{wrapfig}
\usepackage{url}
\usepackage{color}
\usepackage{epsfig}
\usepackage{multirow}
\usepackage{colortbl}

\usepackage{graphicx}
\graphicspath{ {imagenes/} }
\usepackage[font=small,justification=centering]{caption}
%\usepackage{subcaption}

\usepackage[left=3cm,right=2.5cm,top=2.5cm,bottom=2.5cm]{geometry}
\usepackage{tikz}
\usetikzlibrary{calc}

\usepackage{fontspec}
\usepackage{times} 
\usepackage{mathptmx}
\usepackage[font={footnotesize,it,bf}]{caption}



\usepackage{datetime}
\newdateformat{monthyeardate}{
  \monthname[\THEMONTH]  \THEYEAR}  \makeatletter
\renewcommand{\monthnamespanish}[1][\month]{%
  \@orgargctr=#1\relax
  \ifcase\@orgargctr
    \PackageError{datetime}{Invalid Month number \the\@orgargctr}{%
      Month numbers should go from 1 to 12}%
    \or Enero%
    \or Febrero%
    \or Marzo%
    \or Abril%
    \or Mayo%
    \or Junio%
    \or Julio%
    \or Agosto%
    \or Septiembre%
    \or Octubre%
    \or Noviembre%
    \or Diciembre%
    \else \PackageError{datetime}{Invalid Month number \the\@orgargctr}{%
      Month numbers should go from 1 to 12}%
  \fi}
\makeatother

\usepackage{titlesec}
\titleformat{\section}
  {\normalfont\fontsize{12}{14}\bfseries}{\thesection .}{1em}{}
\titleformat{\subsection}
  {\normalfont\fontsize{12}{14}\bfseries}{\thesubsection .}{1em}{}

\usepackage[hidelinks]{hyperref} % Para usar hipervinculos
\usepackage[]{verbatim} % Para poder comentar por bloques
\usepackage[]{lipsum} % Para rellenar con texto

\title{De-reverberación del habla a partir de algoritmos de aprendizaje profundo}

\author{Meza, Martin Bernardo}

\begin{document}
%\setlength{\parindent}{12pt}

\setmainfont{Carlito}
%\setsansfont[Scale=MatchLowercase]{Calibri}
\usetikzlibrary{calc}
\makeatletter
\begin{titlepage}
	\begin{tikzpicture}[remember picture, overlay]
  	\draw[line width = 1pt] ($(current page.north west) + (1.5cm,-2cm)$) rectangle ($(current page.south east) + (-1.5cm,2cm)$);
	\end{tikzpicture}
    \begin{center}
	   	   
	   \includegraphics[width=0.6\textwidth]{imagenes/logo.png} 
	          
        \vspace*{0.36cm}
        \Large
        \fontsize{18pt}{18pt}\selectfont\
        \textbf{INGENIERÍA DE SONIDO}
               
        \vspace{4.5cm}
        \fontsize{22pt}{22pt}\selectfont\
        \textbf{\@title \dag}
        
        \vspace{3.5cm}
        \fontsize{18pt}{18pt}\selectfont\
        \textbf{Autor: \@author }
        
        \fontsize{16pt}{16pt}\selectfont\
        \textbf{Tutores: }
        
        \vspace{3.5cm}
        
	  \normalsize(\dag) \textbf{Tesis para optar por el titulo de ingeniero/a de Sonido}.
        
        \vspace{0.8cm}
        
        
        \large
        \monthyeardate\today
        
    \end{center}
\end{titlepage}
\makeatother
\newpage
\setcounter{page}{1}
\setcounter{section}{0}
\setmainfont{Carlito}
\begin{center}
    \textbf{Plan de Tesis}
\end{center}
\section[Introducción]{CAPÍTULO 1:$\ \ \ \ $INTRODUCCIÓN} 

\subsection[Fundamentación]{FUNDAMENTACIÓN}
Las tecnologías de procesamiento digital de señales de voz mostraron grandes avances en las últimas décadas, llegando a ocupar un rol importante en nuestro día a día. Las investigaciones realizadas en este campo impulsaron diversas aplicaciones basadas en el análisis de la voz humana \cite{fun1}\cite{fun2}. 
Estas aplicaciones, en mayor o menor medida, deben lidiar con una característica intrínseca a cualquier emisión sonora dentro de un recinto: la reverberación. Esto se debe principalmente a que las señales de voz se obtienen a partir de un transductor que no siempre se encuentra cercano a la fuente que se desea registrar, provocando que la señal registrada contenga la reverberación propia del entorno. Esta reverberación interfiere con la señal de voz, produciendo una reducción en el rendimiento de aquellas aplicaciones que dependen de la integridad de dicha señal, como por ejemplo: 

Reconocimiento del habla \cite{reconocimiento}, verificación del hablante\footnote{Debe distinguirse entre reconocimiento del habla y verificación del hablante. Lo primero refiere a poder distinguir qué palabras fueron dichas, y lo segundo refiere a identificar quién es el que está pronunciando las palabras.} \cite{verificacion} y localización del hablante \cite{localizacion}.

Si bien esta problemática fue abordada desde el enfoque de diversas técnicas de procesamiento de señales, en los últimos años ocurrieron grandes avances producto de la implementación de una tecnología emergente de amplio crecimiento en el ambiente científico: los algoritmos de aprendizaje profundo. La capacidad y robustez de estos métodos a la hora de resolver problemas pertinentes al procesamiento de imágenes y detección de patrones se vio también reflejada en el campo de la dereverberación de habla. Actualmente, los sistemas basados en modelos de aprendizaje profundo representan el estado del arte, tanto en dereverberación del habla como otras tareas de audio, desplazando a enfoques mas clásicos del procesamiento de señales. Sin embargo, aún quedan desafíos por resolver como por ejemplo: la falta de bases de datos masivas de señales acústicas específicas como respuestas al impulso reales, la selección de una representación óptima de las señales que permita explotar sus características intrínsecas, las formas de evaluar y cuantificar el desempeño de los sistemas, entre otros.

Por este motivo, esta investigación pretende realizar un análisis de estas problemáticas desde el punto de vista de la ingeniería de sonido, para comprender las limitaciones de los modelos actualmente utilizados en este campo de estudio, analizar alternativas posibles a la escasez de datos y poder aportar al progreso y mejora del rendimiento de dichos modelos. 

\subsection[Objetivos]{OBJETIVOS}
\subsubsection{Objetivo general}

El objetivo general de esta investigación es implementar un algoritmo de dereverberación de señales de voz a partir del uso de redes neuronales y algoritmos de aprendizaje profundo. 

\subsubsection{Objetivos específicos}

\begin{itemize}
    \item Realizar una revisión de las técnicas utilizadas para resolver el problema de dereverberación.
    \item Diseñar e implementar una arquitectura de red neuronal para dereverberación de señales de voz en lenguaje Python.
    \item Analizar técnicas de pre y post procesamiento de datos, estudiando el impacto que tienen en el rendimiento del algoritmo.
    \item Optimizar el sistema propuesto, y evaluar los resultados obtenidos de manera objetiva.
    \item Estudiar y analizar técnicas de generación y aumentación de datos, evaluando su impacto en el desempeño del sistema implementado. 
\end{itemize}

\subsection[Estructura de la Investigación]{ESTRUCTURA DE LA INVESTIGACIÓN}
En el capítulo 2 se presenta el estado del arte referido a las técnicas de dereverberación de señales del habla. 
En el capítulo 3 se detalla el marco teórico necesario para el seguimiento y comprensión de este trabajo. En este se abordan tres temáticas principales: la representación de señales de audio en el dominio espectral mediante la transformada de tiempo corto de Fourier, el concepto de reverberación y su relación con la respuesta al impulso, y por último la aplicación de redes neuronales convolucionales y algoritmos de aprendizaje junto con las principales técnicas de procesamiento.  
En el capítulo 4 se especifica de manera detallada la metodología seguida a lo largo de este trabajo, y se brinda toda la información necesaria para replicar los experimentos realizados. 
En el capítulo 5 se presentan los resultados de los experimentos y se hace un análisis crítico de los mismos. 
En el capítulo 6 se exponen las conclusiones generales del trabajo, y por último en el capítulo 7 se proponen líneas futuras de investigación relacionadas con el presente trabajo. 
\section[Estado del Arte]{CAPÍTULO 2:$\ \ \ \ $ESTADO DEL ARTE} 

En los últimos años se ha registrado un marcado desarrollo y progreso en el campo del procesamiento de señales del habla. En este campo, la dereverberación ocupa un rol crucial, debido al impacto negativo que genera la presencia de reverberación en muchas aplicaciones. 

Los primeros enfoques que apuntaron a resolver el problema de la dereverberación se orientaron al modelado o registro de las respuestas al impulso y la estimación de filtros inversos a partir de estas \cite{filtros_inv}. Como el efecto de la reverberación en una señal se puede pensar como el resultado de una convolución entre una señal anecoica y una respuesta al impulso, este enfoque apunta a estimar la respuesta al impulso con el fin de poder generar un filtro inverso que permita realizar una deconvolución de la señal para poder revertir el efecto de la respuesta del recinto, recuperando la señal en su estado anecoico. Sin embargo, esta metodología presenta varios inconvenientes, como el hecho de considerar que las respuestas al impulso provienen de sistemas lineales e invariantes en el tiempo, lo cual no siempre se cumple en la práctica \cite{LTI}, o bien, el hecho de que la respuesta no siempre pueda ser deducida de manera directa y deba ser estimada. 


También surgieron trabajos enfocados en modelar matemáticamente la señal de habla anecoica  \cite{rabiner}. Algunos de estos trabajos consistían en estimar la señal de habla mediante predicción lineal y calcular el residuo, el cual contiene información sobre la reverberación en la señal. Esta señal de residuo se utilizó para estimar filtros variantes en el tiempo que, al aplicarse a la señal de habla, lograban eliminar parte de la reverberación \cite{LPresiduo}. Otro enfoque consistió en utilizar múltiples transductores y aplicar técnicas de factorización matricial, como la descomposición en valores singulares (SVD), sobre las señales captadas \cite{multichannel}. Algunas características propias de las señales de habla, como la estructura armónica \cite{armonica} y el índice de modulación  \cite{mod}, fueron explotadas para eliminar los efectos de la reverberación. 

Posteriormente, se aplicó la idea de la sustracción espectral \cite{spect_subtrac} \cite{spect_subtrac2} que básicamente consiste en la estimación del espectro de potencia generado por la reverberación a partir de modelos estadísticos. En 2006, Wang et. al. aplicaron este enfoque combinado con el de la estimación de filtros inversos logrando presentar avances importantes en la efectividad de los algoritmos \cite{two_stage}.
 
 
El uso de máscaras binarias ideales en el dominio temporal-frecuencial para extraer las señales buscadas \cite{binarymask} es un enfoque muy utilizado, el cual tiene su origen en el campo del análisis computacional de escenas auditivas \cite{ASA}. Las máscaras se definen como ideales ya que su obtención requiere del conocimiento de la señal buscada y de la señal que interfiere. El uso de estas máscaras implica primero realizar una transformación de la señal de entrada de manera de trasladarla al dominio tiempo-frecuencia (por ejemplo, calculando un espectrograma o un cocleograma) y luego asignarle a cada punto del espacio temporal-frecuencial un valor de 1 cuando su energía mayormente pertenece a la señal objetivo, y un valor de 0 en el caso contrario. Roman et. al. \cite{rev_mask} aplicaron este concepto para tratar el problema de dereverberación, donde se busca estimar la máscara binaria ideal tomando como señal objetivo la señal del habla en condiciones anecoicas y como interferencia a la parte reverberante. Para conseguir la dereverberación, este método requiere seleccionar de manera correcta parámetros como el punto desde el cual se distingue la parte temprana y tardía de la reverberación, y el nivel del umbral en base al cual se identifica a un punto específico como parte de la señal deseada o como parte de la interferencia \cite{parametros}. Hazrati et al. \cite{hazrati} propusieron estimar la máscara binaria a partir de un parámetro dependiente de la varianza de la señal, la cual define un umbral adaptativo, obteniendo mejores resultados. 

A partir del año 2007, se comenzaron a aplicar redes neuronales en la tarea de dereverberación. Jin y Wang \cite{MLP} utilizaron perceptrones multicapa para estimar las máscaras binarias necesarias para la separación de la componente reverberante en una señal voz. La red neuronal aprende a estimar máscaras binarias a partir de la representación tiempo-frecuencia de la señal reverberada. Mas adelante, con el avance de los modelos de aprendizaje profundo, esta técnica se iría perfeccionando reflejándose en mejores resultados en la tarea de dereverberación. Los enfoques basados en la estimación de máscaras tuvieron variantes como, por ejemplo, la estimación de máscaras ideales reales \cite{cIRM}, máscaras ideales complejas \cite{IRM} y máscaras sensibles a la fase \cite{GAN}.
  
En 2014 Kun et al. \cite{ezeKun} proponen el uso de redes neuronales profundas entrenadas para que sean capaces de estimar el espectro de señales anecoicas a partir de señales reverberantes. Nuevamente se plantea la búsqueda del filtro inverso que permita la deconvolución de la señal reverberante para obtener su versión anecoica, pero usando redes neuronales para estimar ese filtro. 

Se han explorado una gran cantidad de arquitecturas y tipos de redes neuronales profundas. Las más utilizadas son las redes neuronales convolucionales (CNN), que surgieron del estudio de la corteza visual del cerebro y han sido muy exitosas en algunas tareas visuales complejas \cite{lagartija}.  Las CNN en general trabajan sobre espectrogramas de magnitud, y tienen una estructura de codificador-decodificador \cite{FCN, rir_filtinverso}. Estas son eficientes en términos de parámetros, aunque requieren de una gran cantidad de capas (o profundidad). Esto se debe a que cada capa convolucional modela su entrada de forma local, con un campo receptivo limitado, y es necesario colocar muchas capas en serie para ampliar ese campo receptivo y abarcar la totalidad del espectrograma de entrada. Otras arquitecturas han sido exploradas para dereverberación del habla, como por ejemplo redes recurrentes \cite{RNN}, combinaciones de redes recurrentes y convolucionales \cite{RNN+CNN} y redes generativas adversarias \cite{GAN}.


Gran parte de los trabajos procesan la magnitud del espectrograma, ignorando la información de fase. En estos casos, al realizar la inversión del espectrograma estimado, algunos sistemas \cite{ezeKun} utilizan el algoritmo de Griffin Lim \cite{griffinlim}, el cual permite invertir espectrogramas utilizando solamente su magnitud. Otros trabajos utilizan la fase original de la señal reverberada \cite{CNN, FCN, skip, rir_filtinverso}, lo cual es una solución sencilla aunque subóptima. Por último, en trabajos recientes la información de fase es utilizada por las redes neuronales, ya sea porque trabajan con el espectrograma complejo \cite{cIRM}, o porque modelan directamente la forma de onda \cite{hifiGAN}.
\section{Marco Teórico}


\subsection{Representación temporal y frecuencial de señales}

\subsection{Respuesta al impulso y reverberación}

\subsection{Inteligibilidad y parámetros de calidad de percepción}

\subsection{Redes neuronales y algoritmos de aprendizaje}



\section{Diseño de la Investigación}
Esta investigación busca estudiar el rendimiento de algoritmos de aprendizaje profundo en la tarea de la segregación de las componentes reverberantes de una señal de voz. Para conseguirlo, se propone organizar la metodología en las siguientes etapas:\\

\textit{Etapa 1: Revisión del estado del arte y planteo de objetivos} \\
Lo primero será estudiar detalladamente los documentos precedentes en este campo específico de aplicación. Se debe prestar especial atención a los trabajos mas recientes para poder determinar claramente cuales son los puntos débiles que presentan las metodologías actuales y de ese modo dilucidar de que manera se puede aportar al problema en cuestión y desde que enfoque conviene hacerlo. Con esto, se busca tener una mayor perspectiva sobre el problema para poder determinar con mayor eficiencia los objetivos pretendidos a corto y largo plazo. \\

\textit{Etapa 2: Determinar las herramientas a utilizar} \\
Con los objetivos en mente, el siguiente paso es determinar con que herramientas desarrollar los algoritmos para poder resolver cada paso de la cadena de procesamiento requerida. Existiendo tanta variedad de opciones, se debe analizar detalladamente que elección resulta mas fructífera para la implementación de los procesos pretendidos y para el posterior análisis de dicho proceso. Esto incluye tanto las herramientas de software o frameworks, tanto como los datos o bases de datos a utilizar. \\

\textit{Etapa 3: Análisis y Pre-procesamiento de los datos} \\
Habiendo elegido las herramientas y los datos a utilizar, el siguiente paso sera realizar un análisis exploratorio de dichos datos. Este análisis tiene como finalidad lograr comprender de una mejor manera las características de los datos con los que se va a trabajar, y determinar cuales son los procesamientos previos que se deben aplicar para que estos datos sean útiles a la hora de alimentar el algoritmo de aprendizaje profundo. Probablemente esta etapa sea la de mayor importancia en cuanto a que es el campo en donde se puede realizar un mayor aporte desde el punto de vista de la Ingeniería de Sonido. En esta etapa se deberá determinar con que datos alimentar el algoritmo y de que manera (o que características) serán estos introducidos a la cadena de procesamiento. Los resultados de este análisis y las decisiones que se deriven de estos marcaran el rumbo de la investigación. \\

\textit{Etapa 4: Modelado y puesta a punto de la/las estructuras de aprendizaje profundo a utilizar, junto con su posterior entrenamiento}\\
En esta etapa se realiza la implementación del algoritmo de aprendizaje profundo. Consiste tanto en armar la estructura de red elegida previamente, como también el resto de la cadena de procesamiento (previo y posterior a la red). Se deben definir las métricas a utilizar para cuantificar el rendimiento de los algoritmos, las funciones de  optimización para la propagación del error, y otros parámetros relativos al algoritmo. \\

\textit{Etapa 5: Análisis de los resultados de cada modelo, validación y comparación con los modelos existentes.}\\
En esta etapa de la investigación se busca hacer un contraste con el trabajo realizado y los antecedentes existentes, de manera de valorar los resultados obtenidos y ponerlos en contexto. Se deben elegir adecuadamente los parámetros a comparar y deben considerarse todas las decisiones tomadas a lo largo del procesamiento para poder abordar a conclusiones representativas sobre el aporte de la investigación al campo de estudio. \\

\textit{Etapa 6: Armado de interfaz gráfica}\\
Por último, la etapa restante consiste en el armado de una herramienta que permita visualizar el funcionamiento del algoritmo propuesto. Esta etapa también contempla el armado de una función global que permita utilizar un modelo entrenado para ser insertado en una cadena de procesamiento posterior, para ser utilizado en otras aplicaciones como el de-noising. \\


\subsection{Cronograma}

En el siguiente diagrama de Grantt se propone el cronograma de actividades para la realización de esta investigación. 

\begin{table}[H]
\begin{tabular}{|l|l|l|l|l|l|l|l|l|l|l|l|l|}
\cline{1-1}
Actividad / Quincenas                                                                             & 1 & 2 & 3 & 4 & 5 & 6 & 7 & 8 & 9 & 10 & 11 & 12 \\ \hline
\begin{tabular}[c]{@{}l@{}}Revision de estado del arte  y \\ \\ planteo de objetivos\end{tabular} & X & X & X &   &   &   &   &   &   &    &    &    \\ \hline
Determinación de herramientas                                                                     &   & X & X &   &   &   &   &   &   &    &    &    \\ \hline
Análisis y pre-procesamiento de datos                                                             &   &   &   & X & X & X & X &   &   &    &    &    \\ \hline
Modelado de las estructuras de aprendizaje profundo                                               &   &   &   &   &   &   & X & X &   &    &    &    \\ \hline
Análisis y validación de resultados                                                               &   &   &   &   &   &   &   &   & X & X  & X  &    \\ \hline
Desarrollo de Interfaz de visualización de datos                                                  &   &   &   &   &   &   &   &   &   &    & X  & X  \\ \hline
\end{tabular}
\end{table}
%\section{Validación}
%\section{Análisis de los resultados}
%\section{Conclusiones}
A lo largo de este trabajo se implementó un algoritmo de aprendizaje profundo para la tarea de dereverberación de señales de habla. Se utilizó una estructura de tipo autoencoder, la cual se entrenó para estimar máscaras de amplitud que al aplicarse sobre un espectro reverberado generen un espectro dereverberado, que luego de procesarse pueda desembocar en información de audio dereverberado. Particularmente se estudió la influencia del manejo de datos para el entrenamiento de este algoritmo, manipulando y generando respuestas al impulso de diversas maneras. 
%\section{Lineas futuras de investigación}

\newpage

\printbibliography[heading=bibintoc,title={Bibliografia}]

%\appendix
%\section{}
%\section{Proceso de aumentación de $T_{60}$}
asdasd




\end{document}
